\documentclass[12pt]{article}
\usepackage[T1]{fontenc}
\usepackage[utf8]{inputenc}
\usepackage{lmodern}
\usepackage{textcomp}
\usepackage{enumitem}
\usepackage{geometry}
\usepackage{graphicx}
\usepackage{amsmath, amssymb}
\geometry{margin=1in}
\begin{document}
\begin{center}
Geography - Graduate Level Assessment 18 \
Total Marks: 40 \
Answer all questions.
\end{center}

\section*{Multiple Choice (10 marks)}
\begin{enumerate}[label=\arabic*.]
\item As of 2019, about what percentage of Turks say God plays an important role in their life?
\begin{enumerate}[label=(\alph*)]
    \item 59\%
    \item 69\%
    \item 79\%
    \item 89\%
\end{enumerate}
\item Which of the following is not an element of the redistribution-with-growth policy approach?
\begin{enumerate}[label=(\alph*)]
    \item minimum wage legislation
    \item land reform
    \item progressive taxation
    \item increased access to education
\end{enumerate}
\item How many children between the ages of 5-14 worked globally as of 2012?
\begin{enumerate}[label=(\alph*)]
    \item 5 million
    \item 30 million
    \item 150 million
    \item 500 million
\end{enumerate}
\item As of 2015, agriculture made up about what percentage of total US GDP?
\begin{enumerate}[label=(\alph*)]
    \item 1\%
    \item 3\%
    \item 9\%
    \item 20\%
\end{enumerate}
\item As of 2019, about what percentage of people from the United States says homosexuality should be accepted by society?
\begin{enumerate}[label=(\alph*)]
    \item 52\%
    \item 62\%
    \item 72\%
    \item 82\%
\end{enumerate}
\end{enumerate}

\section*{True / False (10 marks)}
\textit{Answer True or False}
\begin{enumerate}[label=\arabic*.]
\setcounter{enumi}{5}
\item True or False: Edson Luiz Sampel, a Brazilian expert in canon law, says that canon law is contained in the genesis of various institutes of civil law, such as the law in continental a different answer countries.
\item True or False: 2013 Economics Nobel prize winner Robert J. Shiller said that rising inequality in the United States and elsewhere is the most important problem. Increasing inequality harms economic growth. High and persistent unemployment, in which inequality increases, has a negative effect on subsequent long-run economic growth.
\item True or False: After the end of hostilities, a period of growth and expansion started for the city. In 1920 a stagecoach bus line was established joining Montevideo with the newly formed settlement of Unión and the first natural gas street lights were inaugurated.
\item True or False: Gautama was now determined to complete his spiritual quest. At the age of 35, he famously sat in meditation under a Ficus religiosa tree now called the Bodhi Tree in the town of Bodh Gaya and vowed not to rise before achieving enlightenment.
\item True or False: The Boston Globe wrote that "if there's one school that can lay claim to educating the nation's top national leaders over the past three decades, it's Yale." Yale alumni were represented on the Democratic or Republican ticket in a different answer U.
\end{enumerate}

\section*{Long Form Response (20 marks)}
\textit{Answer each question with a clear and complete explanation.}
\begin{enumerate}[label=\arabic*.]
\setcounter{enumi}{10}
\item Read the following passage and answer the question: Sir John Call argued the advantages of Norfolk Island in that it was uninhabited and that New Zealand flax grew there. In 1786 the British government included Norfolk Island as an auxiliary settlement, as proposed by John Call, in its plan for colonisation of New South Wales. The decision to settle Norfolk Island was taken due to Empress Catherine II of Russia's decision to restrict sales of hemp. Practically all the hemp and flax required by the Royal Navy for cordage and sailcloth was imported from Russia. Question: In including Norfolk Island as an auxiliary settlement, what Australian state did the British government plan to colonise?
\item Read the following passage and answer the question: Macao: The event was held in Macau on May 3. It was the first time that the Olympic torch had traveled to Macau. A ceremony was held at Macau Fisherman's Wharf. Afterward, the torch traveled through Macau, passing by a number of landmarks including A-Ma Temple, Macau Tower, Ponte Governador Nobre de Carvalho, Ponte de Sai Van, Macau Cultural Centre, Macau Stadium and then back to the Fisherman's Wharf for the closing ceremony. Parts of the route near Ruins of St. Paul's and Taipa was shortened due to large crowds of supporters blocking narrow streets. A total of 120 torchbearers participated in this event including casino tycoon Stanley Ho. Leong Hong Man and Leong Heng Teng were the first and last torchbearer in the relay respectively. An article published on Macao Daily News criticized t... Question: Who was the first torchbearer in Macao?
\end{enumerate}

\vfill
\noindent\textit{End of Paper}
\end{document}